\documentclass{amsart}
\usepackage{amsmath,amssymb,amsthm,hyperref}

\theoremstyle{plain}
\newtheorem{theorem}{Theorem}
\newtheorem{lemma}{Lemma}
\theoremstyle{definition}
\newtheorem{definition}{Definition}
\newtheorem{remark}{Remark}

\newcommand{\xA}{x^{*}_{A}}
\newcommand{\xB}{x^{*}_{B}}

\begin{document}

\title{An abstract Gronwall comparison principle via monotone Picard operators}
\maketitle

\section{Set-up and the conditions identified}

We work in the following abstract framework, which is the standard one for
``abstract Gronwall lemmas'' (comparison principles for Picard operators in
ordered spaces).

\begin{definition}[Ordered metric space with closed order]\label{def:oms}
A triple $(X,\le,d)$ is an \emph{ordered metric space} if $(X,d)$ is a metric
space, $(X,\le)$ is a partially ordered set, and the order is
\emph{$d$-closed}: whenever $u_n\to u$ and $v_n\to v$ in $(X,d)$ and
$u_n\le v_n$ for all $n$, then $u\le v$. (Equivalently, the set
$\{(u,v):u\le v\}$ is closed in $X\times X$.)
\end{definition}

\begin{definition}[Picard operator]\label{def:picard}
A self-map $T:X\to X$ on the metric space $(X,d)$ is a \emph{Picard operator}
if $T$ has a unique fixed point $x^{*}_{T}$ and, for \emph{every} $x\in X$, the
sequence of Picard iterates $T^{n}(x)$ converges to $x^{*}_{T}$ in $(X,d)$:
\[
   T^{0}(x):=x,\qquad T^{n+1}(x):=T\bigl(T^{n}(x)\bigr),\qquad
   \lim_{n\to\infty}T^{n}(x)=x^{*}_{T}.
\]
\end{definition}

\begin{remark}\label{rem:examples}
Definition~\ref{def:picard} is exactly the Picard (global asymptotic
stability) property. It holds in many concrete situations, e.g.:
\begin{itemize}
\item $(X,d)$ complete and $T$ a Banach contraction (Banach fixed point
      theorem): then $T$ is a Picard operator;
\item more generally $T$ a $\varphi$-contraction, a Matkowski contraction, or
      any weakly Picard operator with a unique fixed point;
\item $(X,d)$ finite (discrete metric) and every $T$-orbit eventually
      absorbed at a unique fixed point.
\end{itemize}
We do not need a specific contraction estimate; only the orbit-convergence in
Definition~\ref{def:picard} is used.
\end{remark}

A map $T:X\to X$ is \emph{monotone (increasing)} if $u\le v$ implies
$T(u)\le T(v)$. We say $A\le B$ \emph{pointwise} if $A(x)\le B(x)$ for every
$x\in X$.

\medskip

The additional structural and analytic conditions we identify, and under which
the abstract Gronwall comparison principle holds, are the following.

\begin{theorem}[Abstract Gronwall comparison principle]\label{thm:main}
Let $(X,\le,d)$ be an ordered metric space with $d$-closed order
\textup{(Definition~\ref{def:oms})}, and let $A,B:X\to X$ satisfy:
\begin{enumerate}
\item[\textup{(H1)}] $A$ is a Picard operator with unique fixed point $\xA$;
\item[\textup{(H2)}] $B$ is a Picard operator with unique fixed point $\xB$;
\item[\textup{(H3)}] $A\le B$ pointwise, i.e.\ $A(x)\le B(x)$ for all $x\in X$;
\item[\textup{(H4)}] $A$ is monotone increasing.
\end{enumerate}
Then the abstract Gronwall comparison principle holds:
\[
   \forall x\in X,\qquad x\le A(x)\ \Longrightarrow\ x\le \xB .
\]
\end{theorem}

Note that hypotheses (H1)--(H2) refine the bare assumptions of the problem
(unique fixed points) into the Picard property, while (H3) is the dominance
assumption from the problem; the single extra structural assumption is the
monotonicity (H4) of the \emph{dominated} operator $A$ together with the
$d$-closedness of the order. Monotonicity of $B$ is \emph{not} required. The
necessity of (H4) is discussed in Section~\ref{sec:sharp}.

\section{Proof of Theorem~\ref{thm:main}}

The proof rests on two lemmas. Throughout, (H1)--(H4) are assumed.

\begin{lemma}[Every $A$-sub-fixed point lies below $\xA$]\label{lem:sub}
For every $x\in X$, if $x\le A(x)$ then $x\le \xA$.
\end{lemma}

\begin{proof}
Fix $x$ with $x\le A(x)$. We first show by induction on $n\ge 0$ that
\begin{equation}\label{eq:monchain}
   A^{n}(x)\le A^{n+1}(x).
\end{equation}
For $n=0$ this is the hypothesis $x=A^{0}(x)\le A^{1}(x)=A(x)$. Assume
\eqref{eq:monchain} for some $n$, i.e.\ $A^{n}(x)\le A^{n+1}(x)$. Applying the
monotone map $A$ (H4) to this inequality gives
$A\bigl(A^{n}(x)\bigr)\le A\bigl(A^{n+1}(x)\bigr)$, that is
$A^{n+1}(x)\le A^{n+2}(x)$, which is \eqref{eq:monchain} for $n+1$. This
completes the induction.

Next we show, again by induction on $n$, that
\begin{equation}\label{eq:xbelow}
   x\le A^{n}(x)\qquad\text{for all }n\ge 0 .
\end{equation}
For $n=0$, $x\le x$ holds by reflexivity. If $x\le A^{n}(x)$, then combining
with \eqref{eq:monchain} (which gives $A^{n}(x)\le A^{n+1}(x)$) and
transitivity of $\le$ yields $x\le A^{n+1}(x)$. This proves
\eqref{eq:xbelow}.

By (H1), $A$ is a Picard operator, so $A^{n}(x)\to\xA$ in $(X,d)$ as
$n\to\infty$. Consider the two convergent sequences
\[
   u_n:=x\ \ (\text{constant}, \ u_n\to x),\qquad
   v_n:=A^{n}(x)\to \xA .
\]
By \eqref{eq:xbelow} we have $u_n=x\le A^{n}(x)=v_n$ for all $n$. Since the
order is $d$-closed (Definition~\ref{def:oms}), passing to the limit gives
$x\le \xA$.
\end{proof}

\begin{lemma}[The fixed point of $A$ lies below that of $B$]\label{lem:fp}
$\xA\le \xB$.
\end{lemma}

\begin{proof}
The crucial step is the pointwise comparison of the two iteration semigroups,
which we prove by induction on $n\ge 1$:
\begin{equation}\label{eq:semicomp}
   A^{n}(x)\le B^{n}(x)\qquad\text{for all }x\in X .
\end{equation}

\emph{Base case $n=1$.} This is exactly (H3): $A(x)\le B(x)$ for all $x$.

\emph{Inductive step.} Assume \eqref{eq:semicomp} holds for some $n\ge 1$ and
all $x\in X$. Fix $x\in X$ and write $y:=A^{n}(x)$ and $z:=B^{n}(x)$, so that
the inductive hypothesis says $y\le z$. Then:
\[
   A^{n+1}(x)=A(y)\ \overset{(\mathrm{H4})}{\le}\ A(z)\ \overset{(\mathrm{H3})}{\le}\ B(z)=B^{n+1}(x).
\]
Here the first inequality uses that $A$ is monotone (H4) applied to $y\le z$,
and the second uses the pointwise dominance $A(z)\le B(z)$ (H3). By
transitivity, $A^{n+1}(x)\le B^{n+1}(x)$, completing the induction and proving
\eqref{eq:semicomp}.

Now apply \eqref{eq:semicomp} with the particular choice $x=\xA$. Then for
every $n\ge 1$,
\begin{equation}\label{eq:atstar}
   A^{n}(\xA)\le B^{n}(\xA).
\end{equation}
Since $\xA$ is the fixed point of $A$, we have $A^{n}(\xA)=\xA$ for all $n$, so
the left-hand sequence is the constant $\xA$, converging to $\xA$. By (H2),
$B$ is a Picard operator, so the right-hand sequence converges:
$B^{n}(\xA)\to\xB$. Applying the $d$-closedness of the order to the
inequalities \eqref{eq:atstar} between the convergent sequences
$u_n:=A^{n}(\xA)=\xA\to\xA$ and $v_n:=B^{n}(\xA)\to\xB$ yields
$\xA\le \xB$.
\end{proof}

\begin{proof}[Proof of Theorem~\ref{thm:main}]
Let $x\in X$ with $x\le A(x)$. By Lemma~\ref{lem:sub}, $x\le\xA$. By
Lemma~\ref{lem:fp}, $\xA\le\xB$. By transitivity of $\le$, $x\le\xB$. Since
$x$ was an arbitrary $A$-sub-fixed point, the comparison principle holds.
\end{proof}

\section{Sharpness of the hypotheses}\label{sec:sharp}

We record that the structural hypotheses cannot simply be dropped. In both
examples below $X$ is finite and carries the discrete metric, for which the
order is automatically $d$-closed (every convergent sequence is eventually
constant), so Definitions~\ref{def:oms}--\ref{def:picard} apply and the only
hypothesis being violated is the one named.

\subsection*{Monotonicity (H4) is essential.}
Take the three-element poset $X=\{0,1,2\}$ with order generated by $0\le 2$
(and reflexivity), so that $1$ is incomparable to both $0$ and $2$. Define
\[
   A=B:\quad 0\mapsto 2,\quad 1\mapsto 1,\quad 2\mapsto 1 .
\]
Then $A=B$ has the unique fixed point $1$ (indeed $A(0)=2\ne0$, $A(2)=1\ne2$),
and every orbit reaches $1$ in finitely many steps
($0\mapsto2\mapsto1\mapsto1$, $2\mapsto1$), so both $A$ and $B$ are Picard
operators with $\xA=\xB=1$. Trivially $A\le B$ pointwise. However $A$ is
\emph{not} monotone: $0\le 2$ but $A(0)=2\not\le 1=A(2)$. Here $x=0$ satisfies
$0\le 2=A(0)$, i.e.\ $0$ is an $A$-sub-fixed point, yet $0\not\le 1=\xB$ because
$0$ and $1$ are incomparable. Thus the conclusion fails when (H4) is removed.

\subsection*{Dominance (H3) is essential.}
Keep $X=\{0,1,2\}$, but now with order generated by $0\le 1$ (so $2$ is
isolated). Let
\[
   A:\ 0\mapsto0,\ 1\mapsto0,\ 2\mapsto0,\qquad
   B:\ 0\mapsto1,\ 1\mapsto2,\ 2\mapsto2 .
\]
Then $A$ is monotone ($0\le1\Rightarrow A(0)=0\le0=A(1)$), with unique fixed
point $\xA=0$ and all orbits absorbed at $0$, so $A$ is a monotone Picard
operator. The map $B$ has unique fixed point $\xB=2$ and orbits
$0\mapsto1\mapsto2\mapsto2$, $1\mapsto2$, so $B$ is a Picard operator. The
element $x=0$ satisfies $0\le0=A(0)$. But (H3) \emph{fails}: at the point $2$
we have $A(2)=0$ and $B(2)=2$, which are incomparable, so $A(2)\le B(2)$ does
not hold. Correspondingly the conclusion fails: $0\not\le 2=\xB$. This shows the
pointwise dominance (H3) cannot be weakened to hold only on, say, the orbit of
$\xA$.

\begin{remark}[On the role of $B$'s monotonicity]
Theorem~\ref{thm:main} requires monotonicity of $A$ only. This is genuinely
weaker than the frequently-quoted ``both operators increasing'' hypothesis: the
induction \eqref{eq:semicomp} compares the $A$- and $B$-iterates by pushing the
\emph{outer} map through the monotonicity of $A$ and then using dominance, so
$B$ is never required to preserve order. In particular, the conclusion
$\xA\le\xB$ is obtained without assuming that $B$-sub-fixed points lie below
$\xB$ (a statement which is false for non-monotone $B$). Of course, assuming
both $A$ and $B$ monotone is a (stronger) sufficient condition under which
Theorem~\ref{thm:main} likewise applies.
\end{remark}

\section{Specialization: the classical Gronwall situation}

The framework recovers the classical Gronwall lemma pattern. Let $(X,d)$ be a
complete metric space carrying a $d$-closed order $\le$, and suppose $A$ is a
\emph{monotone Banach contraction} and $B$ a Banach contraction with $A\le B$
pointwise. By the Banach fixed point theorem, $A$ and $B$ are Picard operators
(Remark~\ref{rem:examples}), so (H1), (H2) hold; (H3), (H4) are the standing
assumptions. Hence Theorem~\ref{thm:main} applies: every $x$ with $x\le A(x)$
satisfies $x\le\xB$. In the prototypical analytic instance $X=C([0,T])$ with the
pointwise order and the sup metric, $A u(t)=u_0+\int_0^t f(s,u(s))\,ds$ and
$B u(t)=u_0+\int_0^t g(s,u(s))\,ds$ with $f\le g$ and $f$ nondecreasing in its
second argument, this is precisely the differential/integral Gronwall
comparison: a subsolution of the $f$-equation is dominated by the solution of
the $g$-equation.

\end{document}
